\chapter{Vaginismus}
\index{Vaginismus}

\medskip
\noindent\textit{Educational reference; always seek professional advice for individual care.}

\section*{Definition and Overview}
Vaginismus is a condition in which the muscles around the outer third of the vagina contract involuntarily, making any attempt at vaginal penetration difficult, painful, or impossible. The word comes from the Latin \textit{vagina}, reflecting the fact that the problem lies in the pelvic floor muscles rather than in any physical blockage. It is classified as a female sexual pain disorder, and in the past was sometimes referred to as a "psychosexual" condition, but it is best understood as a complex problem in which physical, psychological, and behavioural factors combine. The involuntary muscle spasm may occur during attempts at sexual intercourse, the insertion of a tampon, a gynaecological examination, or a cervical smear test. Importantly, vaginismus is not something a person does deliberately or can simply "relax out of"; it is an automatic reflex, rather like the way the eyelid closes if something approaches the eye. The condition is highly treatable, and with appropriate education, physiotherapy, and gradual desensitisation techniques, most people affected can achieve comfortable penetration and satisfying sexual activity. This chapter provides an overview of vaginismus, its causes, how it is diagnosed, and the treatments that are available.

\section*{Epidemiology}
Accurate figures for vaginismus are difficult to obtain because many affected people never seek help, and studies use different definitions. Estimates of the lifetime prevalence of vaginismus in women of reproductive age generally range from about 1 to 7 percent, although some community surveys report higher figures. It is diagnosed across all ages, ethnicities, and educational backgrounds, and can occur at any point in life, including in women who have previously had comfortable intercourse. Vaginismus is frequently under-recognised: affected women may avoid gynaecological examinations, smear tests, and tampon use for years, and may present only when a relationship or a desire to conceive makes the problem unavoidable. It is also common for vaginismus to coexist with anxiety disorders, a history of sexual abuse, or other sexual pain problems. Because the condition is surrounded by embarrassment and silence, the true number of affected women is likely to be considerably higher than the number who receive a diagnosis.

\section*{Aetiology and Pathophysiology}
Vaginismus is best understood as a learned, automatic reflex in which the pelvic floor muscles, particularly the pubococcygeus and other superficial muscles of the perineum, contract in anticipation of penetration. This contraction is a protective response, and it is maintained by a cycle of fear and pain: the anticipation of pain triggers the spasm, the spasm makes penetration painful, and the pain reinforces the fear.

Contributing factors include:
\begin{itemize}
    \item \textbf{Fear of pain:} past painful experiences with intercourse, tampons, or medical examinations can "teach" the body to anticipate discomfort
    \item \textbf{History of trauma or sexual abuse:} prior physical or sexual assault is an important risk factor, although most women with vaginismus have no such history
    \item \textbf{Anxiety and perfectionism:} women who are generally anxious, or who hold rigid beliefs about sex, are over-represented
    \item \textbf{Negative early learning:} restrictive or negative messages about sex during childhood
    \item \textbf{Physical factors:} conditions that make penetration painful, such as vulval vestibulitis, thrush, dryness, or pelvic inflammatory disease, can initiate the reflex
    \item \textbf{Relationship factors:} pressure to have sex, poor communication, or fear of pregnancy can contribute
\end{itemize}

Importantly, in many women no clear precipitating factor can be identified. Whatever the trigger, the result is the same: a reflex spasm of the pelvic floor that prevents comfortable penetration.

\section*{Clinical Features}
The hallmark of vaginismus is the involuntary contraction of the pelvic floor muscles during any attempt at penetration. The features a woman may notice include:

\begin{itemize}
    \item \textbf{Inability to have intercourse:} attempts at penetration either fail entirely or are extremely painful ("like hitting a wall")
    \item \textbf{Pain with tampon insertion:} difficulty or inability to insert or remove tampons
    \item \textbf{Difficulty with medical examinations:} severe discomfort, tightening, or an inability to tolerate a speculum or a digital examination
    \item \textbf{Pelvic and thigh muscle tension:} the legs may draw together, the back may arch, and the hips may lift off the bed during attempts
    \item \textbf{Pain described as burning, tearing, or stinging} at the vaginal entrance rather than deeper inside
    \item \textbf{Avoidance behaviour:} many women avoid sexual activity, smear tests, and even discussions about the problem, and this can strain relationships
    \item \textbf{Secondary problems:} anxiety, low self-esteem, guilt, and relationship conflict often develop as the problem persists
\end{itemize}

Vaginismus is described as primary when penetration has never been possible, and secondary when it develops after a period of comfortable penetration, for example after childbirth, surgery, or a traumatic event.

\section*{Differential Diagnosis}
Because the symptoms of vaginismus overlap with other causes of pelvic pain and penetration difficulty, a careful assessment is needed to distinguish them. Many of these conditions can coexist, and all should be considered.

\begin{itemize}
    \item \textbf{Vulvodynia or vestibulodynia:} burning pain at the entrance of the vagina without a visible cause, which often coexists with vaginismus
    \item \textbf{Infection:} thrush, bacterial vaginosis, or sexually transmitted infections causing vulval irritation and pain
    \item \textbf{Hormonal dryness:} atrophic vaginitis after menopause, breastfeeding, or use of hormonal contraception
    \item \textbf{Structural problems:} a congenitally narrow or absent vaginal opening, or scarring of the hymen
    \item \textbf{Endometriosis or pelvic inflammatory disease:} causes of deep pelvic pain that can make penetration uncomfortable
    \item \textbf{Dyspareunia from other causes:} a full sexual-pain assessment is required to determine whether the pain is primarily at the entrance or deeper
    \item \textbf{Apasmofilia is not a recognised medical diagnosis} and should not be confused with vaginismus
\end{itemize}

The key distinction is that in vaginismus the pain is caused by involuntary muscle contraction rather than by identifiable tissue damage or disease.

\section*{When to Seek Medical Care}
Many women live with vaginismus for years, assuming the problem is "in their head" or that they are unusual. In fact, it is a recognised, treatable condition, and it is always worth raising with a general practitioner, sexual health clinic, gynaecologist, or women's health physiotherapist.

Seek medical care if:
\begin{itemize}
    \item Penetration is consistently painful or impossible and it is affecting your quality of life
    \item You want to conceive and penetration is preventing intercourse
    \item You have never been able to use a tampon or have a smear test
    \item The problem started after an injury, childbirth, surgery, or a traumatic event
    \item Pain is accompanied by bleeding, unusual discharge, or a burning sensation that suggests infection
\end{itemize}

There is no such thing as "too minor" a reason: if penetration causes distress, professional help is appropriate. No medical emergency is associated with vaginismus itself, but if attempts at intercourse cause severe, persistent pain or bleeding, an assessment is warranted.

\section*{Diagnosis and Investigations}
The diagnosis of vaginismus is based primarily on the history, and a diagnosis is usually made with confidence in the consultation. Sensitive questioning will cover the timing of the problem, whether it is primary or secondary, the situations that trigger it, any history of trauma, and the impact on relationships and well-being. A physical examination is often offered, but it is always conducted at the woman's pace, with her consent, and can be deferred, modified, or conducted partly through her own self-exploration. During the examination the clinician will look for the characteristic spasm, which may be visible as the muscles tightening as soon as the area is approached. Unless there are warning symptoms such as bleeding, discharge, or suspicion of infection, no laboratory tests are required, and imaging is rarely needed. If an underlying physical cause such as infection or vulval disease is suspected, swabs and a referral to a gynaecologist may be arranged. Part of the diagnosis is educational: explaining the reflex nature of the problem and that it is common and treatable is itself therapeutic.

\section*{Management}
The treatment of vaginismus is highly effective and focuses on breaking the cycle of fear, spasm, and pain. A multidisciplinary approach, combining education, physiotherapy, and psychological support, gives the best results.

\begin{itemize}
    \item \textbf{Education and reassurance:} understanding that vaginismus is a reflex, not a character flaw, reduces shame and is the foundation of treatment
    \item \textbf{Pelvic floor physiotherapy:} a women's health physiotherapist teaches conscious relaxation of the pelvic floor, breathing techniques, and graded exercises
    \item \textbf{Dilator therapy:} using smooth, graduated vaginal dilators, combined with lubrication and relaxation techniques, to gradually desensitise the area and gain control over the muscles; this is done at the woman's own pace and can take weeks to months
    \item \textbf{Cognitive behavioural therapy:} addressing unhelpful thoughts about pain, sex, and worthiness, and gradually exposing the person to feared situations
    \item \textbf{Sex therapy or counselling:} helping the woman and, where appropriate, her partner communicate about the problem and rebuild a positive sexual relationship
    \item \textbf{Botulinum toxin injections:} in a small number of resistant cases, injections into the pelvic floor muscles may help break the spasm, allowing other therapy to progress
    \item \textbf{Treating any underlying physical problem:} such as infection or dryness, which can perpetuate the reflex
    \item \textbf{Patience and pacing:} success depends on going slowly, avoiding pain, and celebrating small steps
\end{itemize}

The woman is always in control of the pace. A common misconception is that "just trying harder" or drinking alcohol will help; in practice, gentleness and graded exposure work far better.

\section*{Complications}
Vaginismus itself is not dangerous, but untreated it can have significant consequences for quality of life.

\begin{itemize}
    \item Inability to conceive through intercourse, and the distress this can cause for couples trying for a family
    \item Avoidance of essential health care, including cervical smear tests, which may delay detection of cervical changes
    \item Relationship strain, guilt, and reduced intimacy with a partner
    \item Low self-esteem, shame, and symptoms of anxiety or depression
    \item Secondary anorgasmia or reduced sexual desire as the avoidance and distress escalate
    \item Self-treatment attempts, including the use of objects or excessive force, which can cause injury
\end{itemize}

None of these complications are inevitable, and prompt treatment prevents them from becoming entrenched.

\section*{Prognosis and Outcomes}
The outlook for vaginismus is excellent. With education, physiotherapy, and graduated dilator work, the great majority of women achieve comfortable penetration, and treatment success rates reported in clinical studies are high, typically above 80 to 90 percent. Relapses can occur, particularly at times of stress, childbirth, or gynaecological procedures, but the skills learned in treatment usually allow a woman to manage these without losing progress. Successful treatment restores the ability to have smear tests, use tampons, and enjoy sexual activity, and it also improves confidence and relationship satisfaction. The key messages are that vaginismus is common, is nobody's fault, and responds well to the right kind of help.

\section*{Prevention and Self-Care}
\begin{itemize}
    \item Talk openly about sexual pain: it is a legitimate health topic, not a source of shame
    \item Practise pelvic floor relaxation and breathing exercises regularly
    \item Use plenty of water-based lubricant during any penetration attempt
    \item Progress gradually: never force penetration, and stop at the first sign of pain
    \item Use dilators or graded finger insertion as part of a structured, relaxed home programme
    \item Communicate honestly with a partner about what feels comfortable and what does not
    \item Attend cervical screening using a smaller speculum, or request a "self-test" HPV option, if internal examinations are difficult
    \item Seek help early, rather than avoiding the issue, and consider involving a partner in therapy
\end{itemize}

\section*{Sources and Further Reading}
The following organisations provide reliable information about vaginismus and sexual pain for women, partners, and clinicians.

\begin{itemize}
    \item NHS. \textit{Vaginismus: symptoms, causes, and treatment.} https://www.nhs.uk/
    \item Mayo Clinic. \textit{Women's health: pain with intercourse and vaginismus.} https://www.mayoclinic.org/
    \item MSD Manual. \textit{Vaginismus: consumer overview of female sexual pain disorders.} https://www.msdmanuals.com/
    \item MedlinePlus. \textit{Vaginismus health information.} https://medlineplus.gov/
    \item Jean Hailes for Women's Health. \textit{Sexual pain and vaginismus information for Australian women.} https://www.jeanhailes.org.au/
    \item Family Planning Australia. \textit{Sexual and reproductive health information and services.} https://www.fpnsw.org.au/
\end{itemize}
